\section{Introduction}
Digital recruitment has significantly increased the number of applications received for a single job posting. Online job portals allow applicants to apply across locations with very low effort, and organizations often receive hundreds or thousands of resumes for a single opening. Human recruiters therefore spend substantial time performing repetitive first-stage screening. This stage is also vulnerable to inconsistency because reviewers may give excessive weight to formatting, university names, familiar employers, or visually prominent sections rather than the actual match between candidate capabilities and role requirements \cite{b3}, \cite{b10}.

Applicant Tracking Systems (ATS) were introduced to reduce this workload, but many traditional systems operate primarily through rigid keyword matching. A job description requiring ``Python'', ``REST'', and ``PostgreSQL'' may cause the system to search for these exact terms in candidate resumes. This approach is computationally inexpensive, but it creates false negatives when a candidate uses alternate wording, abbreviations, or semantically related phrasing. It is also vulnerable to keyword stuffing, where candidates insert repeated job-description terms into resumes in order to inflate match scores.

Recent recruitment systems increasingly use transformer-based semantic embeddings and large language models. These methods can map related expressions such as ``RESTful APIs'' and ``REST'' into nearby semantic spaces, reducing lexical brittleness \cite{b6}, \cite{b12}. However, they introduce two practical challenges. First, they increase latency, memory consumption, and deployment cost, particularly when processing large batches. Second, they often behave as black-box scoring layers. A recruiter may receive a high or low match score without a clear, auditable explanation of which job requirements were satisfied, missing, or over-weighted. In hiring contexts, such opacity is not merely inconvenient; it can reduce trust and make candidate rejection decisions difficult to justify \cite{b8}, \cite{b14}.

This paper presents a practical middle-ground architecture for first-pass resume screening. The proposed framework combines deterministic skill extraction, switchable TF-IDF or MiniLM context similarity, bounded weighted scoring, optical character recognition fallback, and rule-based explainability. Rather than claiming complete semantic understanding in all operating modes, it is designed around three constraints common in small and medium-sized enterprise recruitment: speed, low deployment cost, and auditability. Each score is converted into a plain-language explanation that lists matched skills, missing skills, detected soft skills, experience fit, and bonus skills.

The main contributions of this work are:
\begin{itemize}
    \item an end-to-end OCR-aware resume screening pipeline that handles text, PDF, and image resumes using digital extraction followed by OCR fallback;
    \item a bounded weighted scoring model that combines technical skills, soft skills, experience, selectable contextual similarity, and capped extra-skill bonuses;
    \item an explainable AI layer that converts score components into recruiter-readable justifications;
    \item a benchmark-driven comparison of the lightweight Regex+TF-IDF design against semantic embedding alternatives with respect to robustness, efficiency, and interpretability.
\end{itemize}
